% =============================================================================
% Chapter 1: Introduction — ML Systems Lecture Slides (Gold Standard)
% =============================================================================
% This deck is the reference template for all other chapters.
% Expert-reviewed: Patterson x2, Dean, Education x2, Gemini.
%
% IMAGES NEEDED (generate SVGs, convert to PDF):
%   - dam-taxonomy.pdf        - iron-law-visual.pdf
%   - sw1-vs-sw2.pdf          - five-percent-problem.pdf
%   - verification-gap.pdf    - ai-eras-timeline.pdf
%   - compute-growth.pdf      - silent-failure.pdf
%   - degradation-curve.pdf   - deployment-spectrum.pdf
%   - amdahls-pipeline.pdf    - vit-vs-resnet.pdf
%   - resnet-bottleneck.pdf   - monitoring-dashboard.pdf
%   - data-center.jpg         - nvidia-a100.png
% =============================================================================
\documentclass[aspectratio=169, 12pt]{beamer}
\usepackage{../../assets/beamerthememlsys}

\mlsyssetup{
  volume   = {Volume I},
  chapter  = {Chapter 1},
  logo = {../../assets/img/logo-mlsysbook.png},
  instlogo = {../../assets/img/logo-harvard.png},
  chaptertitle = {Introduction},
}

% --- Fonts ---
\usepackage{fontspec}
\setsansfont{Helvetica Neue}[
  BoldFont={Helvetica Neue Bold},
  ItalicFont={Helvetica Neue Italic},
  BoldItalicFont={Helvetica Neue Bold Italic},
]
\setmonofont{JetBrains Mono}[Scale=0.85]

% --- Packages ---
\usepackage{booktabs}
\usepackage{amsmath}

% --- Image paths ---
\graphicspath{
  {images/}
}

% --- Chapter-specific macros ---
\newcommand{\DAM}{D\raisebox{0.08em}{\tiny$\bullet$}A\raisebox{0.08em}{\tiny$\bullet$}M}

% --- Helper: safe image include (compiles even if image missing) ---
% Note: \IfFileExists does NOT respect \graphicspath, so we check images/ explicitly
\newcommand{\safeimg}[2][width=\textwidth,keepaspectratio]{%
  \IfFileExists{images/#2}{\includegraphics[#1]{#2}}{%
    \IfFileExists{#2}{\includegraphics[#1]{#2}}{%
      \fbox{\parbox[c][2.5cm][c]{0.85\linewidth}{\centering\footnotesize\textcolor{midgray}{[Missing image]}}}%
    }%
  }%
}

% --- Section count for navigation (must match actual \section{} count) ---
\setcounter{mlsystotalsections}{8}

\title{Introduction}
\author{Vijay Janapa Reddi}
\institute{Harvard University}
\date{}

\begin{document}

% =============================================================================
% TITLE SLIDE
% =============================================================================
\mlsystitle{Introduction}{The Physics of AI Engineering}{cover_introduction.png}

% =============================================================================
% LEARNING OBJECTIVES
% =============================================================================
\begin{frame}{Learning Objectives}
\note{
% -- NARRATE: Read through objectives aloud. ``This chapter establishes the
% analytical vocabulary for the entire course. By the end, you will have three
% diagnostic tools---D-A-M, the Iron Law, and the Degradation Equation---that
% we will use in every single chapter.''
% Ask: ``How many of you have trained a model but never thought about the
% hardware it runs on?''
%
% -- FLEX: [CORE] Sets the learning contract for the chapter.
% IF SHORT: Display objectives without reading each one; ask the hand-raise
% question to establish engagement.
}

\small
\begin{enumerate}
  \item Explain the \textbf{AI Triad} and apply the \textbf{\DAM{} taxonomy}
  \item Trace AI's evolution and explain the \textbf{Bitter Lesson}
  \item Distinguish ML systems from traditional software
  \item Apply the \textbf{Iron Law of ML Systems} to decompose performance
  \item Identify how \textbf{Lighthouse Models} stress-test bottlenecks
  \item Apply the \textbf{Degradation Equation} to reason about drift
  \item Define \textbf{AI Engineering} and the \textbf{efficiency framework}
  \item Apply the \textbf{five-pillar framework} to organize ML engineering
\end{enumerate}

\end{frame}

\begin{frame}{Visual Language}
\note{
% -- NARRATE: Walk through each card. ``Blue means compute---anytime you see
% blue in a diagram, think GPU ops, matrix multiplies, forward pass. Green is
% data flow---memory, caches, healthy paths. Orange is routing---schedulers,
% load balancers. Red is cost, error, or bottleneck.'' Point to each card as
% you name it.
%
% -- FLEX: [CORE] Show on Day 1 and again briefly on Day 2.
% IF SHORT: Display the slide but do not narrate---students can read the cards.
}

\small
Throughout this course, colors carry meaning:

\vspace{0.3cm}
\begin{columns}[T]
  \begin{column}{0.45\textwidth}
    \begin{mlsyscard}{computestroke}
    \textbf{Blue} --- Compute / Processing\\
    {\footnotesize GPU ops, forward/backward pass, inference}
    \end{mlsyscard}
    \vspace{0.15cm}
    \begin{mlsyscard}{datastroke}
    \textbf{Green} --- Data / Memory\\
    {\footnotesize Data flow, caches, healthy paths}
    \end{mlsyscard}
  \end{column}
  \begin{column}{0.45\textwidth}
    \begin{mlsyscard}{routingstroke}
    \textbf{Orange} --- Routing / Scheduling\\
    {\footnotesize Load balancers, batch windows}
    \end{mlsyscard}
    \vspace{0.15cm}
    \begin{mlsyscard}{errorstroke}
    \textbf{Red} --- Error / Cost / Bottleneck\\
    {\footnotesize Loss, decode phase, waste}
    \end{mlsyscard}
  \end{column}
\end{columns}
\end{frame}


% =============================================================================
\section{The AI Moment}
% =============================================================================

\begin{frame}{The AI Moment}
\note{
% -- LINK: Learning objectives promised students would understand why ML systems
% differ from traditional software. This slide grounds that claim in everyday
% experience.
%
% -- NARRATE: ``Every one of these systems runs inference millions of times a
% day on real hardware. Voice assistants do speech-to-text-to-response in under
% 200 ms. Social media ranks billions of posts per second. Lending decisions
% evaluate creditworthiness under uncertainty. Autonomous driving fuses camera,
% LiDAR, and radar in real time.'' Point to each bullet.
% Ask: ``Which of these did you interact with today?''
%
% -- ENGAGE: ``Pick one of these systems. What physical constraint do you think
% limits its performance?'' Give 30 seconds, cold-call one student.
% Expected: latency for voice, compute for ranking, power for autonomous.
%
% -- WARN: Students think AI = the model. Correct framing: each example is a
% full system---data pipeline, hardware, serving infrastructure, monitoring---not
% just a neural network.
%
% -- FLEX: [CORE] Opens the chapter with visceral connection to students' lives.
% IF SHORT: Skip the cold-call; just ask the ``which did you interact with''
% question as a show of hands.
}

\footnotesize
\begin{columns}[T]
  \begin{column}{0.55\textwidth}
    AI is not a future possibility --- it is present reality:
    \begin{itemize}\setlength\itemsep{0pt}
      \item \textbf{Voice assistants}: speech $\to$ text $\to$ response
      \item \textbf{Social media}: ranking billions of posts in ms
      \item \textbf{Lending}: creditworthiness under uncertainty
      \item \textbf{Autonomous driving}: fusing sensor modalities
    \end{itemize}
    \vspace{0.05cm}
    \begin{mlsyscard}{crimson}
      These systems \alert{make decisions under uncertainty}.
    \end{mlsyscard}
  \end{column}
  \begin{column}{0.42\textwidth}
    \safeimg{data-center.jpg}
  \end{column}
\end{columns}

\end{frame}

\begin{frame}{The Dual Mandate}
\note{
% -- LINK: The AI Moment showed systems making decisions under uncertainty.
% This slide names the two simultaneous constraints every such system faces.
%
% -- NARRATE: ``Every ML system juggles two things at once: statistical
% uncertainty---predictions are probabilistic, accuracy degrades silently---and
% physical constraints---memory, power, latency, speed of light. Traditional
% software crashes loudly when it fails. ML systems degrade silently. Look at
% the table: a bug in traditional software is a code error that crashes. A bug
% in ML is a data error that degrades predictions without any stack trace.''
%
% -- ENGAGE: ``Give me one example of a silent failure in an ML system you have
% used.'' Cold-call one student. Expected: recommendation quality declining,
% autocomplete getting worse, ad relevance dropping.
%
% -- WARN: Students equate ML with ``better software.'' Correct framing: ML
% introduces a fundamentally new failure mode---silent probabilistic
% degradation---that traditional testing cannot catch.
%
% -- FLEX: [CORE] This is the chapter's thesis statement.
% IF AHEAD: ``Which column of the table is harder to debug, and why?''
}

\small
\begin{columns}[T]
  \begin{column}{0.48\textwidth}
    Every ML system must simultaneously manage:

    \vspace{0.2cm}
    \begin{mlsyscard}{computestroke}
    \textbf{Statistical uncertainty}\\
    {\footnotesize Predictions are probabilistic. Accuracy degrades silently.}
    \end{mlsyscard}

    \vspace{0.15cm}

    \begin{mlsyscard}{routingstroke}
    \textbf{Physical constraints}\\
    {\footnotesize Memory, power budgets, latency, speed of light.}
    \end{mlsyscard}
  \end{column}
  \begin{column}{0.48\textwidth}
    \vspace{0.3cm}
    \renewcommand{\arraystretch}{1.25}
    {\footnotesize
    \begin{tabular}{@{}lll@{}}
      \toprule
      & \textbf{Traditional} & \textbf{ML System} \\
      \midrule
      \textbf{Bug}   & Code error   & Data error \\
      \textbf{Failure} & Crash (loud) & Silent degradation \\
      \textbf{Debug} & Trace code   & Inspect data \\
      \bottomrule
    \end{tabular}
    }
  \end{column}
\end{columns}

\end{frame}

% =============================================================================
\section{The Paradigm Shift}
% =============================================================================

\begin{frame}{Software 1.0 vs.\ Software 2.0}
\note{
% -- LINK: The Dual Mandate showed that ML systems face silent degradation.
% This slide explains the root cause: in Software 2.0, data IS the source code.
%
% -- NARRATE: ``Karpathy's framing makes the shift concrete. In SW 1.0, the
% programmer writes explicit logic. In SW 2.0, the programmer curates the
% dataset, and SGD is the compiler. Look at the table: source code is now
% training data. The compiler is SGD. Logic is learned, not written. Failure
% is metric decay, not a crash. Debugging moves upstream---from code to data.''
% Point to diagram: ``The arrow from data to model is the key insight.''
%
% -- ENGAGE: ``If data is source code, what does version control look like?''
% Give 20 seconds. Cold-call one student.
% Expected: dataset versioning tools like DVC. Deepen: ``What happens if
% someone deletes a row of training data without tracking it?''
%
% -- WARN: Students think ``data is important'' but still treat the model as
% the primary artifact. Correct framing: in SW 2.0, the dataset is the
% deliverable. The model is derived from it, like a compiled binary.
%
% -- FLEX: [CORE] Foundational conceptual shift for the course.
% IF SHORT: Skip the engage question; narrate the table for 90 seconds.
}

\small
\begin{columns}[T]
  \begin{column}{0.45\textwidth}
    \safeimg[width=\textwidth,height=5.5cm,keepaspectratio]{sw1-vs-sw2.pdf}
  \end{column}
  \begin{column}{0.52\textwidth}
    \renewcommand{\arraystretch}{1.2}
    {\footnotesize
    \begin{tabular}{@{}lll@{}}
      \toprule
      & \textbf{SW 1.0} & \textbf{SW 2.0} \\
      \midrule
      Source & C++, Python & Training data \\
      Compiler & GCC, LLVM & SGD \\
      Logic & Explicit & Learned \\
      Failure & Crash & Metric decay \\
      Debug & Trace path & Inspect data \\
      \bottomrule
    \end{tabular}
    }
    \vspace{0.2cm}

    \mlsysconcept{Key insight:}{Data is source code.}

    \vspace{0.15cm}
    Debugging moves \emph{upstream} --- from code to data.
  \end{column}
\end{columns}
\mlsyscite{Karpathy, ``Software 2.0,'' 2017}

\end{frame}

\begin{frame}{The 5\% Problem}
\note{
% -- LINK: SW 2.0 showed that data is source code. This slide reveals the
% scale of the non-model infrastructure---95\% of a production ML system.
%
% -- NARRATE: ``Sculley et al.\ called this `hidden technical debt.' Point to
% the tiny center box in the diagram: that is the ML model code---about 5
% percent. Everything around it---data pipelines, feature stores, serving
% infrastructure, monitoring, configuration---is what this course teaches.
% The model is the tip of the iceberg.''
%
% -- ENGAGE: ``Look at the diagram. What is the biggest surrounding box?''
% Give 15 seconds. Expected answer: data collection or configuration.
% Both are valid---the point is that neither is the model.
%
% -- WARN: Students equate ``ML'' with ``the model.'' Correct framing: the
% model is 5\%; the other 95\% is systems engineering that determines whether
% the model ever reaches a user.
% IF STUCK: ``Think about what happens after you call model.predict()---who
% routes the request, monitors latency, detects drift?''
%
% -- FLEX: [CORE] Foundational mental model for the course.
% IF SHORT: Skip the question; just narrate the diagram for 90 seconds.
}

\small
\begin{columns}[T]
  \begin{column}{0.50\textwidth}
    \safeimg[width=\textwidth,height=5.5cm,keepaspectratio]{five-percent-problem.pdf}
  \end{column}
  \begin{column}{0.46\textwidth}
    \begin{mlsyscard}{errorstroke}
    \textbf{Hidden Technical Debt}\\[0.1cm]
    {\footnotesize In production ML, the model is $\approx$5\% of the code. The other 95\%:}
    \begin{itemize}
      \item {\footnotesize Data pipelines \& feature engineering}
      \item {\footnotesize Serving infrastructure}
      \item {\footnotesize Monitoring \& configuration}
    \end{itemize}
    \end{mlsyscard}

    \vspace{0.15cm}
    \textbf{``ML is easy; ML Systems are hard.''}
  \end{column}
\end{columns}
\mlsyscite{Sculley et al., NeurIPS 2015}

\end{frame}

\begin{frame}{The Verification Gap}
\note{
% -- LINK: The 5\% Problem showed the iceberg of systems infrastructure.
% This slide explains WHY that infrastructure must include monitoring: the
% input space is too vast to test exhaustively.
%
% -- NARRATE: ``A 224-by-224 RGB image has 256 to the power of 150,528
% possible configurations---a number with 362,000 digits. ImageNet tests only
% 50,000 of them. The gap between what we test and what exists is effectively
% infinite. Traditional software can enumerate edge cases. ML systems cannot.
% The engineering consequence: we trade guaranteed correctness for statistical
% monitoring.''
%
% -- ENGAGE: ``If you cannot test all inputs, how do you know your model is
% working?'' Give 20 seconds. Cold-call one student.
% Expected: monitoring production metrics. Deepen: ``What if you do not have
% ground truth labels in real time?''
%
% -- WARN: Students think more test data solves the verification gap. Correct
% framing: 50,000 vs.\ $256^{150,528}$ is not a sampling problem---it is a
% fundamental impossibility. No finite test set can cover the space.
%
% -- FLEX: [OPTIONAL] Can be summarized in one sentence if short on time.
% IF SHORT: ``The input space is infinite; we cannot test our way to
% correctness. Monitoring is non-negotiable.'' Move on.
}

\small
\begin{columns}[T]
  \begin{column}{0.48\textwidth}
    \safeimg[width=\textwidth,height=5.5cm,keepaspectratio]{verification-gap.pdf}
  \end{column}
  \begin{column}{0.48\textwidth}
    A $224\times224$ RGB image has:\\[0.1cm]
    $256^{150{,}528}$ possible configurations\\
    {\footnotesize (a number with \textbf{362,000+ digits})}

    \vspace{0.2cm}
    ImageNet tests only \textbf{50,000} of them.

    \vspace{0.15cm}
    $$\text{Gap} \approx \text{Total Input Space}$$

    \vspace{0.1cm}
    \mlsysalert{Consequence:}{Trade \emph{guaranteed correctness} for \emph{statistical monitoring}.}
  \end{column}
\end{columns}

\end{frame}

% =============================================================================
\section{AI Evolution}
% =============================================================================

\begin{frame}{Seven Decades of Bottlenecks}
\note{
% -- LINK: The Verification Gap showed that ML systems face fundamentally
% different engineering challenges. This slide traces how we got here---seven
% decades of systems bottlenecks, not algorithmic dead ends.
%
% -- NARRATE: ``Walk through four eras. Symbolic AI hit the knowledge
% bottleneck---you cannot hand-code enough rules. Expert systems hit the
% brittleness bottleneck---they broke outside narrow domains. Statistical ML
% hit the feature engineering bottleneck---humans cannot design features for
% every domain. Deep learning hit the compute and data bottleneck---but that
% one we can throw hardware at. The pattern: every AI winter was a systems
% failure, not an algorithmic dead end.''
%
% -- ENGAGE: ``Why did expert systems fail in the 1980s? Was it bad algorithms
% or bad infrastructure?'' Give 15 seconds. Expected: infrastructure---they
% could not scale knowledge acquisition.
%
% -- WARN: Students think AI winters were caused by ``wrong algorithms.''
% Correct framing: each era's algorithms were sound; the systems could not
% deliver the data, compute, or integration required.
%
% -- FLEX: [CORE] Historical context is essential for the Bitter Lesson.
% IF SHORT: Focus on the bottleneck pattern across eras, not individual era
% details. ``Every winter was a systems failure.'' (60 seconds)
}

% --- Layout: FULL-WIDTH IMAGE + caption ---
\centering
\safeimg[width=0.92\textwidth,height=4.8cm,keepaspectratio]{ai-eras-timeline.pdf}

\vspace{0.15cm}
\mlsysinsight{Pattern:}{Every AI winter was a \emph{systems} failure, not an algorithmic dead end.}

\end{frame}

\begin{frame}{The Bitter Lesson}
\note{
% -- LINK: The timeline showed that systems bottlenecks---not algorithms---caused
% every AI winter. Sutton's Bitter Lesson formalizes this pattern.
%
% -- NARRATE: ``Sutton wrote this in 2019. The lesson: general methods that
% leverage computation always win over hand-crafted cleverness. It is `bitter'
% because our intuition says encode expertise---chess engines with opening
% books, NLP with grammar rules. Reality says: scale computation. AlphaGo
% beat hand-tuned heuristics with brute-force search plus learning. GPT beat
% hand-crafted NLP with massive compute on massive data. The implication for
% systems engineers: if scale wins, then the systems that deliver scale are
% the bottleneck.''
%
% -- ENGAGE: ``Why is this lesson bitter? What does it say about human
% intuition?'' Give 20 seconds. Cold-call one student.
% Expected: our intuition misleads us---we want to build clever solutions,
% but brute-force scaling wins.
%
% -- WARN: Students interpret the Bitter Lesson as ``algorithms do not matter.''
% Correct framing: algorithms matter, but only those that can exploit scale.
% The next slide (Efficiency Counter-Narrative) shows algorithmic efficiency
% also compounds.
%
% -- FLEX: [CORE] Philosophical anchor for why systems engineering matters.
% IF AHEAD: ``If scale always wins, why not just buy more GPUs?'' (Answer:
% cost, power, memory wall---the physical constraints from the Course Overview.)
}

\small
\begin{columns}[T]
  \begin{column}{0.55\textwidth}
    Richard Sutton (2019):

    \vspace{0.1cm}
    \begin{mlsyscard}{crimson}
    {\footnotesize ``General methods that leverage computation are
    ultimately the most effective --- and by a large margin.''}
    \end{mlsyscard}

    \vspace{0.1cm}
    The lesson is ``bitter'' because:
    \begin{itemize}\setlength\itemsep{0pt}
      \item Our intuition says \emph{encode expertise}
      \item Reality says \emph{scale computation}
    \end{itemize}

    \vspace{0.05cm}
    {\footnotesize\mlsysconcept{Implication:}{Requires \textbf{systems engineering}.}}
  \end{column}
  \begin{column}{0.42\textwidth}
    \safeimg[width=\textwidth,height=5.5cm,keepaspectratio]{compute-growth.pdf}

    {\footnotesize\textcolor{midgray}{AI compute doubles every 3.4 mo --- 7$\times$ faster than Moore's Law.}}
  \end{column}
\end{columns}
\mlsyscite{Sutton, ``The Bitter Lesson,'' 2019}

\end{frame}

\begin{frame}{The Efficiency Counter-Narrative}
\note{
% -- LINK: The Bitter Lesson said scale wins. This slide shows that algorithmic
% efficiency also compounds---44$\times$ less compute for the same accuracy.
%
% -- NARRATE: ``Two forces co-evolve. Brute-force scaling: AI compute grew 7
% orders of magnitude. Algorithmic efficiency: 44$\times$ compute reduction for
% equivalent accuracy. Both are real. The systems engineer must exploit both.
% But here is the catch: ViT achieves 1--2\% higher accuracy than ResNet-50 but
% requires 4$\times$ the FLOPs and 3$\times$ the memory bandwidth. Better
% accuracy does not mean better system. A model that violates your latency SLA
% has zero production utility.''
%
% -- ENGAGE: ``Which is growing faster---compute demand or algorithmic
% efficiency?'' Give 15 seconds. Expected: compute demand (7$\times$ faster
% than Moore's Law vs.\ 44$\times$ over a decade).
%
% -- WARN: Students assume newer architectures are always better systems choices.
% Correct framing: ViT is a better algorithm but a worse system for latency-
% constrained deployments. ``Better accuracy $\neq$ better system.''
%
% -- FLEX: [OPTIONAL] Safe to compress into one sentence if running short.
% IF SHORT: ``Algorithmic efficiency compounds too, but better accuracy does
% not mean better system---ViT needs 4$\times$ the FLOPs of ResNet.''
}

\small
\begin{columns}[T]
  \begin{column}{0.48\textwidth}
    \safeimg[width=\textwidth,height=5.5cm,keepaspectratio]{vit-vs-resnet.pdf}
  \end{column}
  \begin{column}{0.48\textwidth}
    \textbf{Two forces co-evolve:}

    \vspace{0.15cm}
    \textcolor{computestroke}{\textbf{Brute-force scaling}}\\
    {\footnotesize AI compute: 7 orders of magnitude growth}

    \vspace{0.1cm}
    \textcolor{datastroke}{\textbf{Algorithmic efficiency}}\\
    {\footnotesize 44$\times$ compute reduction for equivalent accuracy}

    \vspace{0.2cm}
    \begin{mlsyscard}{datastroke}
    {\footnotesize ViT: 1--2\% higher accuracy than ResNet-50, but 4$\times$ FLOPs and 3$\times$ memory bandwidth. \textbf{Better accuracy $\neq$ better system.}}
    \end{mlsyscard}
  \end{column}
\end{columns}
\mlsyscite{Hernandez \& Brown, ``Measuring Algorithmic Efficiency,'' 2020}

\end{frame}

\begin{frame}{Return on Compute (RoC)}
\note{
% -- LINK: The Efficiency Counter-Narrative showed that algorithmic efficiency
% compounds. RoC formalizes this as an economic invariant.
%
% -- NARRATE: ``Return on Compute measures how much useful work you extract per
% dollar of compute. The formula: RoC = (Accuracy Gain) / (Compute Cost).
% A model that achieves 1\% accuracy gain for \$100 of compute has an RoC of
% 0.01\%/\$. A model that achieves the same 1\% for \$10,000 has 100$\times$
% worse RoC. The Bitter Lesson says scale wins---but RoC asks: at what price?
% This is the economic discipline that prevents blind scaling.''
%
% -- ENGAGE: ``If GPT-4 cost \$100M to train and improved accuracy by 5\%
% over GPT-3, what is its RoC? Is that good or bad?'' Give 15 seconds.
% Expected: 0.05/100M = 5$\times 10^{-9}$\%/\$. Compare to fine-tuning
% a smaller model: 2\% gain for \$1,000 = 0.002\%/\$.
%
% -- WARN: Students equate ``more compute = better.'' Correct framing: RoC
% forces you to ask whether the marginal accuracy gain justifies the cost.
%
% -- FLEX: [OPTIONAL] Safe to compress if running short.
% IF SHORT: State the formula and the GPT-4 example. (60 seconds)
}

\small
\begin{columns}[T]
  \begin{column}{0.50\textwidth}
    \textbf{Return on Compute:}

    \vspace{0.15cm}
    $$\text{RoC} = \frac{\Delta\text{Accuracy}}{\text{Compute Cost}}$$

    \vspace{0.1cm}
    {\footnotesize The economic invariant that governs scaling decisions.}

    \vspace{0.15cm}
    \mlsysconcept{Key insight:}{The Bitter Lesson says \emph{scale wins}. RoC asks \emph{at what price?}}
  \end{column}
  \begin{column}{0.46\textwidth}
    \vspace{0.2cm}
    \renewcommand{\arraystretch}{1.15}
    {\footnotesize
    \begin{tabular}{@{}lrr@{}}
      \toprule
      \textbf{Approach} & \textbf{Gain} & \textbf{Cost} \\
      \midrule
      GPT-4 pretrain & +5\% & \$100M \\
      Fine-tune small & +2\% & \$1K \\
      Distill + quant & +0\% & \$500 \\
      \bottomrule
    \end{tabular}
    }

    \vspace{0.15cm}
    \begin{mlsyscard}{routingstroke}
    {\footnotesize Fine-tuning: 100,000$\times$ better RoC than pretraining from scratch.}
    \end{mlsyscard}
  \end{column}
\end{columns}

\end{frame}

% --- ACTIVE LEARNING 1: Predict ---
\begin{frame}{Predict: What Defines an ML System?}
\note{
% -- NARRATE: ``Before I show you the answer, write down 3 components that
% define an ML system. Not just `data and model'---think about what makes it
% a \emph{system}.'' This primes the D-A-M taxonomy reveal on the next slide.
%
% -- ENGAGE: Give students 60 seconds to write. Then turn-and-talk for 30
% seconds. Cold-call 2--3 students to share.
% Expected answers: data, model, hardware, serving, monitoring.
% Surprise answer: the interactions between these components.
%
% -- FLEX: [CORE] This prediction exercise is essential scaffolding for the
% D-A-M taxonomy.
% IF SHORT: Cut turn-and-talk; go straight to cold-call after 60 seconds.
}

\centering
\vspace{1.2cm}
{\Large\bfseries Think--Write--Share}

\vspace{0.6cm}
{\large If ML systems are not ``traditional software + a model,''\\
what are the essential components?}

\vspace{0.6cm}
{\normalsize Write down 3 components. \textcolor{midgray}{(60 seconds)}}

\vspace{0.6cm}
{\small\textcolor{lightgray}{Turn to a neighbor and compare your answers.}}

\end{frame}

% =============================================================================
\section{The DAM Taxonomy}
% =============================================================================

\begin{frame}{The D\raisebox{0.04em}{\tiny$\bullet$}A\raisebox{0.04em}{\tiny$\bullet$}M Taxonomy}
\note{
% -- LINK: Students just predicted the essential components of an ML system.
% D-A-M formalizes their intuition into three interdependent axes.
%
% -- NARRATE: ``Reveal after prediction exercise. Three axes, color-coded:
% green = data, blue = algorithm/compute, orange = machine. Point to each in
% the diagram. Data: how much, how fast can we move it---measured in GB and
% GB/s. Algorithm: how many operations, what parallelism pattern---measured in
% FLOPs. Machine: what silicon, what memory hierarchy---measured in FLOP/s.
% The diagnostic question is always the same: which axis is the bottleneck?''
%
% -- ENGAGE: ``Which axis is the bottleneck for training GPT-3?'' Give 15
% seconds. Expected: compute (massive FLOPs). Correct but incomplete---the
% Iron Law will show all three terms matter.
%
% -- WARN: Students treat the three axes as independent knobs. Correct framing:
% they are interdependent---doubling data volume requires proportionally more
% bandwidth or longer training time. Optimizing one shifts pressure to the others.
% IF STUCK: ``Think of a three-legged stool. Shortening one leg does not help
% if another leg was already the shortest.''
%
% -- FLEX: [CORE] First framework---everything builds on this.
% IF AHEAD: ``Can you think of a model where all three axes are equally
% bottlenecked?'' (Answer: rare in practice---one almost always dominates.)
}

\small
\begin{columns}[T]
  \begin{column}{0.48\textwidth}
    \safeimg[width=\textwidth,height=5.5cm,keepaspectratio]{dam-taxonomy.pdf}
  \end{column}
  \begin{column}{0.48\textwidth}
    Every ML system has three interdependent axes:

    \vspace{0.2cm}
    \textcolor{datastroke}{\textbf{Data}} ($D_{\text{vol}}$, $BW$)\\
    {\footnotesize Volume, quality, distribution, bandwidth}

    \vspace{0.15cm}
    \textcolor{computestroke}{\textbf{Algorithm}} ($O$, architecture)\\
    {\footnotesize Operation count, model structure}

    \vspace{0.15cm}
    \textcolor{routingstroke}{\textbf{Machine}} ($R_{\text{peak}}$, memory)\\
    {\footnotesize Hardware throughput, capacity}

    \vspace{0.2cm}
    \mlsysalert{Diagnostic:}{``Which axis is the bottleneck?''}
  \end{column}
\end{columns}

\end{frame}

\begin{frame}{Quick Check}
\note{
% -- NARRATE: ``Quick recall---no peeking. What are the three axes of D-A-M?
% 30 seconds.'' Walk around the room.
%
% -- FLEX: [CORE] Micro-retrieval. Distributed practice improves retention
% (Roediger \& Karpicke, 2006). Never skip retrieval checks.
}

\centering
\Large\bfseries Quick check --- no peeking.\\[0.5cm]
\normalsize What are the three axes of the D$\cdot$A$\cdot$M taxonomy?\\[0.3cm]
{\small 30 seconds --- then we continue.}
\end{frame}



\begin{frame}{D\raisebox{0.04em}{\tiny$\bullet$}A\raisebox{0.04em}{\tiny$\bullet$}M Interdependence}
\note{
% -- LINK: The D-A-M taxonomy named three axes. This slide shows they are
% coupled---optimizing one shifts pressure to the others.
%
% -- NARRATE: ``AlexNet 2012 is the canonical example of D-A-M co-design.
% CNNs existed for decades, but they needed GPU parallelism to become practical.
% The algorithm change (conv layers) aligned with the machine (GPU SIMD).
% Result: error dropped from 26.2\% to 15.3\%---a 42\% relative improvement.
% Not because the algorithm was new, but because the machine could finally
% execute it efficiently. This is D-A-M interdependence in action.''
%
% -- ENGAGE: ``If you switch from ResNet to a Transformer for image
% classification, which other axis feels the most pressure?'' Give 15 seconds.
% Expected: Data (Transformers need more training data) or Machine (more
% memory bandwidth for attention).
%
% -- WARN: Students think they can optimize one axis in isolation. Correct
% framing: changing the algorithm changes the data and hardware requirements.
% AlexNet succeeded because all three axes aligned simultaneously.
%
% -- FLEX: [OPTIONAL] Safe to cover quickly if running short.
% IF SHORT: ``AlexNet won because CNN parallelism matched GPU hardware.
% Axes are coupled.'' (30 seconds) Move on.
}

\footnotesize
\begin{columns}[T]
  \begin{column}{0.55\textwidth}
    \textbf{The axes are coupled:}
    \begin{itemize}\setlength\itemsep{0pt}
      \item \textcolor{computestroke}{Algorithm} change $\to$ different \textcolor{datastroke}{data} needs
      \item \textcolor{datastroke}{Data} change $\to$ different \textcolor{routingstroke}{machine} needed
      \item \textcolor{routingstroke}{Machine} change $\to$ different \textcolor{computestroke}{algorithms} practical
    \end{itemize}

    \vspace{0.05cm}
    {\scriptsize\textbf{AlexNet (2012):} CNNs' parallel ops aligned with GPU hardware. Error: 26.2\% $\to$ 15.3\% --- a 42\% relative improvement.}
  \end{column}
  \begin{column}{0.42\textwidth}
    \safeimg[width=\textwidth,height=4.5cm,keepaspectratio]{nvidia-h100.png}
  \end{column}
\end{columns}

\end{frame}

% =============================================================================
\section{The Iron Law}
% =============================================================================

\mlsysfocus{The Iron Law of ML Systems}{%
$T_{\text{total}} = \underbrace{\dfrac{D_{\text{vol}}}{BW}}_{\text{Data}} +
\underbrace{\dfrac{O}{R_{\text{peak}} \cdot \eta}}_{\text{Compute}} +
\underbrace{L_{\text{lat}}}_{\text{Overhead}}$%
}

\begin{frame}{Iron Law: Three Terms}
\note{
% -- LINK: D-A-M tells you where the bottleneck is. The Iron Law quantifies
% how long each axis takes in seconds.
%
% -- NARRATE: ``Point to each term. Data term: bytes divided by bandwidth
% gives seconds. Compute term: FLOPs divided by peak rate times efficiency
% gives seconds. Overhead: orchestration tax, also seconds. You add three
% times and the slowest one dominates. This is dimensional analysis---if
% your units do not resolve to seconds, the equation is wrong.''
%
% -- ENGAGE: ``For a phone camera app classifying a photo, which term
% dominates?'' Give 20 seconds. Expected: Data term (reading the image from
% memory) or Overhead (framework launch cost). Accept either with reasoning.
%
% -- WARN: Students try to add FLOPs to bytes. Correct framing: every term
% must resolve to seconds before you can compare or add them. Write the units
% under each fraction to verify.
% IF STUCK: ``What are the units of Bytes divided by Bytes-per-second?''
%
% -- FLEX: [CORE] Central equation of the course.
% IF SHORT: State the equation, emphasize ``slowest term dominates,'' move on.
% IF AHEAD: ``What happens to the equation when you pipeline data loading
% with compute?'' (Answer: sum becomes max---the Bottleneck Principle from Ch2.)
}

% --- Layout: FULL-WIDTH diagram + compact legend ---
\centering
\safeimg[width=0.88\textwidth,height=3.8cm,keepaspectratio]{iron-law-visual.pdf}

\vspace{0.15cm}
\small
\begin{columns}[T]
  \begin{column}{0.30\textwidth}
    \centering
    \textcolor{datastroke}{\textbf{Data}} $\frac{D_{\text{vol}}}{BW}$\\[0.05cm]
    {\scriptsize Bytes / (Bytes/s) = \textbf{s}}
  \end{column}
  \begin{column}{0.30\textwidth}
    \centering
    \textcolor{computestroke}{\textbf{Compute}} $\frac{O}{R_p \cdot \eta}$\\[0.05cm]
    {\scriptsize FLOPs / (FLOPs/s) = \textbf{s}}
  \end{column}
  \begin{column}{0.30\textwidth}
    \centering
    \textcolor{errorstroke}{\textbf{Overhead}} $L_{\text{lat}}$\\[0.05cm]
    {\scriptsize Orchestration tax = \textbf{s}}
  \end{column}
\end{columns}
\vspace{0.1cm}
{\small The \textbf{slowest term dominates} end-to-end latency.}

\end{frame}

\begin{frame}{Iron Law: Worked Example}
\note{
% -- LINK: The previous slide defined the Iron Law abstractly. This slide
% applies it to GPT-3 training with real hardware numbers.
%
% -- NARRATE: ``Walk through GPT-3 step by step. Data term: 350 GB divided
% by 2,039 GB/s = 0.17 seconds. Compute term: 3.14 times 10-to-the-14
% FLOPs divided by 989 TFLOPS times 0.5 efficiency = 0.64 seconds. Overhead:
% synchronization and scheduling adds about 0.05 seconds. Total: 0.86 seconds
% per batch. The reveal: compute dominates for LLM training. Upgrading memory
% bandwidth alone---making the data term even faster---yields diminishing
% returns because compute is already 3.7$\times$ slower.''
%
% -- ENGAGE: ``If you upgrade from HBM3 to HBM4 (2$\times$ bandwidth), how
% much does total time improve?'' Give 20 seconds.
% Expected: Data term halves from 0.17 to 0.085 s, saving 0.085 s out of
% 0.86 s = about 10\% improvement. Compute still dominates.
%
% -- WARN: Students assume upgrading any component helps proportionally.
% Correct framing: only upgrading the dominant term yields significant speedup.
% Upgrading the non-dominant term is wasted investment.
%
% -- FLEX: [CORE] First worked example---establishes the analytical pattern.
% IF SHORT: Just show the final comparison: compute at 0.64 s vs.\ data at
% 0.17 s. ``Compute dominates. Case closed.''
}

\small
\textbf{Training one batch of GPT-3 (175B parameters)}

\vspace{0.15cm}
\renewcommand{\arraystretch}{1.2}
{\footnotesize
\begin{tabular}{@{}llll@{}}
  \toprule
  \textbf{Term} & \textbf{Variables} & \textbf{Calculation} & \textbf{Time} \\
  \midrule
  \textcolor{datastroke}{Data}    & $D_{\text{vol}}$=350 GB, $BW$=2,039 GB/s & $350/2{,}039$ & \textbf{0.17 s} \\
  \textcolor{computestroke}{Compute} & $O$=$3.14\!\times\!10^{14}$, $R_p$=989 TF & $3.14\!\times\!10^{14}/(989\!\times\!10^{12}\!\times\!0.5)$ & \textbf{0.64 s} \\
  \textcolor{errorstroke}{Overhead} & $L_{\text{lat}}$ (sync, sched.) & --- & \textbf{0.05 s} \\
  \midrule
  & & $T_{\text{total}} \approx$ & \textbf{0.86 s} \\
  \bottomrule
\end{tabular}
}

\vspace{0.2cm}
\begin{mlsyscard}{computestroke}
\textbf{Diagnosis:} The \textcolor{computestroke}{compute term} dominates for LLM training.
Upgrading memory bandwidth alone yields diminishing returns.
\end{mlsyscard}

\end{frame}

\begin{frame}{Iron Law: ResNet-50 --- A Different Bottleneck}
\note{
% -- LINK: GPT-3 was compute-bound. This slide shows the same equation applied
% to ResNet-50 inference at batch size 1---a completely different bottleneck.
%
% -- NARRATE: ``Same Iron Law, different dominant term. At batch size 1,
% ResNet-50 has an arithmetic intensity of 79 FLOPs per byte. The accelerator
% can compute faster than memory can feed it. The data term dominates---the
% opposite of GPT-3 training. This is why batch size matters: at batch 256,
% compute reuse increases and the compute term catches up.''
%
% -- ENGAGE: ``What happens to the dominant term if you increase batch size
% from 1 to 256?'' Give 15 seconds.
% Expected: Compute term grows (more FLOPs), data term grows less (parameter
% reuse). System shifts from data-bound to compute-bound.
%
% -- WARN: Students assume a model is always bound by the same term. Correct
% framing: the dominant term changes with batch size, hardware, and precision.
% The same model can be data-bound at batch 1 and compute-bound at batch 256.
%
% -- FLEX: [OPTIONAL] Reinforces the Iron Law pattern but can be skipped.
% IF SHORT: Skip this slide and go to the exercise.
% IF AHEAD: ``At what batch size does ResNet-50 cross from data-bound to
% compute-bound on an H100?''
}

\small
\begin{columns}[T]
  \begin{column}{0.48\textwidth}
    \safeimg[width=\textwidth,height=5.5cm,keepaspectratio]{resnet-bottleneck.pdf}
  \end{column}
  \begin{column}{0.48\textwidth}
    \textbf{ResNet-50 inference (batch = 1)}

    \vspace{0.15cm}
    {\footnotesize
    \renewcommand{\arraystretch}{1.15}
    \begin{tabular}{@{}lr@{}}
      \toprule
      \textbf{Metric} & \textbf{Value} \\
      \midrule
      Operations ($O$) & 7.7 GFLOPs \\
      Parameters       & 97.5 MB \\
      Arithmetic Intensity & 79 FLOPs/B \\
      \bottomrule
    \end{tabular}
    }

    \vspace{0.15cm}
    \begin{mlsyscard}{datastroke}
    {\footnotesize At batch 1, the accelerator can compute faster than memory can feed it. \textbf{Data term dominates} --- the opposite of GPT-3 training.}
    \end{mlsyscard}
  \end{column}
\end{columns}

\end{frame}

% --- ACTIVE LEARNING 2: Exercise ---
\begin{frame}{Your Turn: Diagnose the Bottleneck}
\note{
% -- LINK: Students have seen two worked examples (GPT-3 compute-bound,
% ResNet-50 data-bound). Now they apply the Iron Law themselves.
%
% -- NARRATE: ``You have 90 seconds. Calculate all three terms, add them up,
% and tell me which term dominates. Show your units---if they do not resolve
% to seconds, something is wrong.''
%
% -- ENGAGE: PEER INSTRUCTION PROTOCOL:
% 1. Present problem (30s). 2. Individual work (60s).
% 3. If 30--70\% get the right dominant term: pair discussion (90s), re-vote.
% 4. Reveal: Data = 40/100 = 400 ms. Compute = 2e9/(20e9*0.6) = 167 ms.
% Overhead = 5 ms. Total = 572 ms. Data-bound.
% Key follow-up: ``Adding a faster GPU yields 0\% speedup. Why?''
%
% -- WARN: Students will divide FLOPs by peak rate and forget the efficiency
% factor $\eta$. Correct: $R_{\text{peak}} \times \eta = 20 \times 0.6 = 12$
% GFLOPS effective. Without $\eta$, they get 100 ms instead of 167 ms and
% might still identify data as dominant, but the number is wrong.
% IF STUCK: ``Start with the data term. What are the units of 40 MB divided
% by 100 MB/s?''
%
% -- FLEX: [CORE] First hands-on calculation. Do not skip.
% IF SHORT: Skip pair discussion; go straight from individual work to reveal.
}

\small
\begin{columns}[T]
  \begin{column}{0.58\textwidth}
    {\normalsize\bfseries Diagnose the Bottleneck}

    \vspace{0.2cm}
    A mobile inference pipeline has:
    \begin{itemize}
      \item $D_{\text{vol}}$ = 40 MB, $BW$ = 100 MB/s
      \item $O$ = $2\!\times\!10^9$ FLOPs, $R_{\text{peak}}$ = 20 GF, $\eta$ = 0.6
      \item $L_{\text{lat}}$ = 5 ms
    \end{itemize}

    \vspace{0.15cm}
    \textbf{Calculate} $T_{\text{total}}$ and identify the dominant term.

    {\footnotesize\textcolor{midgray}{(90 seconds --- then compare with a neighbor)}}
  \end{column}
  \begin{column}{0.38\textwidth}
    \pause
    \begin{mlsyscard}{datastroke}
    \textbf{Solution:}\\[0.1cm]
    {\footnotesize
    Data: $40/100$ = 400 ms\\
    Compute: $\frac{2\times10^9}{12\times10^9}$ = 167 ms\\
    Overhead: 5 ms\\[0.1cm]
    $T_{\text{total}}$ = \textbf{572 ms}\\[0.1cm]
    \textcolor{datastroke}{\textbf{Data-bound!}}
    }
    \end{mlsyscard}
  \end{column}
\end{columns}

\end{frame}

% =============================================================================
\section{ML vs.\ Software}
% =============================================================================

\begin{frame}{Silent Degradation}
\note{
% -- LINK: The Iron Law measures performance at a point in time. This slide
% asks: what happens to that performance over months?
%
% -- NARRATE: ``This is the most important conceptual slide in the chapter.
% Traditional software crashes loudly---you get a stack trace, a log entry,
% an alert. ML systems degrade silently. Look at the curve: accuracy decays
% 10--15\% over months without a single code change. No crash. No error log.
% No alert unless you built monitoring. The world changed---user behavior
% shifted, data distributions drifted---and the model silently got worse.''
%
% -- ENGAGE: ``How would you detect a 10\% accuracy drop if no one filed a
% bug and no code changed?'' Give 20 seconds. Cold-call one student.
% Expected: monitoring production metrics, comparing to baseline.
% Deepen: ``What if you do not have ground truth labels in real time?''
%
% -- WARN: Students think ``no crash = working correctly.'' Correct framing:
% the absence of errors is not evidence of correctness. ML systems require
% continuous monitoring because their most dangerous failure mode is silence.
%
% -- FLEX: [CORE] Motivates the Degradation Equation on the next slide.
% IF SHORT: Show the diagram, state ``ML fails silently,'' move on (60 seconds).
}

% --- Layout: FULL-WIDTH comparison diagram ---
\centering
\safeimg[width=0.92\textwidth,height=4.5cm,keepaspectratio]{silent-failure.pdf}

\vspace{0.2cm}
\begin{columns}[T]
  \begin{column}{0.45\textwidth}
    \centering\footnotesize
    \textbf{Traditional:} Failure is \emph{loud}\\
    Crashes $\to$ stack traces $\to$ fix
  \end{column}
  \begin{column}{0.45\textwidth}
    \centering\footnotesize
    \textbf{ML:} Failure is \emph{silent}\\
    10--15\% decay over months
  \end{column}
\end{columns}
\vspace{0.1cm}
\mlsysalert{Key:}{No code changed. The \emph{world} changed.}

\end{frame}

\begin{frame}{The Degradation Equation}
\note{
% -- LINK: Silent Degradation showed the phenomenon. This slide formalizes it
% into a measurable equation.
%
% -- NARRATE: ``Accuracy at time t equals initial accuracy minus alpha times
% the distribution distance. Alpha is how sensitive the model is to drift---
% measurable from historical data. Delta measures how far the live data has
% drifted from training data---measurable with PSI or KL divergence. Look at
% the example: a recommendation system starts at 85\% and drops to 79.2\% in
% 6 months. No code changed. No bugs. The engineering response: set a
% retraining trigger when $A(t)$ crosses a threshold.''
%
% -- ENGAGE: ``At what month should this team retrain?'' Give 15 seconds.
% Expected: when accuracy crosses some threshold like 80\%. Deepen: ``How do
% you choose the threshold?'' (Answer: SLA requirements + safety margin.)
%
% -- WARN: Students think drift is rare or only affects ``bad'' models.
% Correct framing: ALL deployed models drift. The question is not whether
% but when. Alpha and Delta are always positive in production.
% IF STUCK: ``Think about a movie recommendation system. Do your tastes
% today match your tastes 6 months ago?''
%
% -- FLEX: [CORE] Third framework---completes the analytical toolkit.
% IF SHORT: Show equation, state the rec-system example, move on.
}

\small
\begin{columns}[T]
  \begin{column}{0.52\textwidth}
    $$A(t) = A_0 - \alpha \cdot \Delta(P_{\text{train}},\; P_{\text{live}}(t))$$

    \vspace{0.1cm}
    {\footnotesize
    \begin{tabular}{@{}ll@{}}
      $A(t)$    & Accuracy at time $t$ \\
      $A_0$     & Accuracy at deployment \\
      $\alpha$  & Sensitivity to drift \\
      $\Delta$  & Distribution distance \\
    \end{tabular}
    }

    \vspace{0.2cm}
    \mlsysconcept{Engineering use:}{Set retraining triggers when $A(t)$ crosses a threshold.}
  \end{column}
  \begin{column}{0.45\textwidth}
    \safeimg[width=\textwidth,height=5.5cm,keepaspectratio]{degradation-curve.pdf}

    {\footnotesize\textcolor{midgray}{Rec.\ system: 85\% $\to$ 79.2\% in 6 months without retraining.}}
  \end{column}
\end{columns}

\end{frame}

\begin{frame}{Defining AI Engineering}
\note{
% -- LINK: Students now have three frameworks (D-A-M, Iron Law, Degradation
% Equation). This slide defines the discipline that uses them.
%
% -- NARRATE: ``AI Engineering is the discipline of designing, deploying, and
% maintaining systems whose outputs are inherently probabilistic to meet
% deterministic reliability targets. Three distinctions from ML Research: (1)
% stochastic outputs vs.\ deterministic SLAs, (2) all three D-A-M axes
% simultaneously, (3) production means continuously shifting distributions.
% In a Kaggle competition, you optimize one metric on a static dataset. In
% production, you optimize accuracy AND latency AND cost on a moving target.''
%
% -- ENGAGE: ``How is this different from what you do in a Kaggle competition?''
% Give 15 seconds. Cold-call one student.
% Expected: single metric vs.\ multi-objective, static vs.\ shifting data.
%
% -- WARN: Students think AI Engineering = ML Research + deployment scripts.
% Correct framing: AI Engineering jointly optimizes all three D-A-M axes under
% continuously shifting distributions---a fundamentally different discipline.
%
% -- FLEX: [CORE] Defines the field the course teaches.
% IF SHORT: Read the definition, show the table, skip the engage question.
}

\small
\begin{mlsyscard}{crimson}
\textbf{AI Engineering:} The discipline of designing, deploying, and maintaining
systems whose outputs are \emph{inherently probabilistic} to meet
\emph{deterministic reliability targets} by satisfying constraints on all three
\DAM{} axes simultaneously in production.
\end{mlsyscard}

\vspace{0.2cm}
\renewcommand{\arraystretch}{1.2}
{\footnotesize
\begin{tabular}{@{}lll@{}}
  \toprule
  & \textbf{ML Research} & \textbf{AI Engineering} \\
  \midrule
  Objective     & Single metric (accuracy)      & Multi-objective (accuracy + latency + cost) \\
  Data          & Static benchmark              & Continuously shifting distribution \\
  Optimization  & Algorithm axis only           & All three \DAM{} axes jointly \\
  Success       & State-of-the-art on test set  & Meets SLA in production \\
  \bottomrule
\end{tabular}
}

\end{frame}

\begin{frame}{The ML System Lifecycle}
\note{
% -- LINK: AI Engineering was defined as maintaining systems under shifting
% distributions. This slide shows the operational shape of that maintenance:
% a continuous loop, not a linear pipeline.
%
% -- NARRATE: ``ML systems require continuous iteration---not a linear handoff.
% Data flows to training, training to deployment, deployment to monitoring,
% and monitoring feeds back to data and retraining. The monitoring loop is the
% structural difference from traditional software. A/B tests in ML take 4--6
% weeks because you need statistical significance on probabilistic outputs.
% Traditional A/B tests take 2--3 days.''
%
% -- ENGAGE: ``Why do ML A/B tests take 4--6 weeks instead of 2--3 days?''
% Give 15 seconds. Expected: need statistical significance on noisy metrics,
% or need enough production data to detect small accuracy differences.
%
% -- WARN: Students think deployment is the end of the project. Correct
% framing: deployment is the beginning of the feedback loop. The model starts
% degrading the moment it goes live.
%
% -- FLEX: [OPTIONAL] Reinforces the loop concept but can be compressed.
% IF SHORT: ``Not a pipeline---a loop. Monitoring feeds back to retraining.''
% Show diagram for 30 seconds, move on.
}

% --- Layout: FULL-WIDTH loop diagram + one-liner ---
\centering
\safeimg[width=0.80\textwidth,height=4.0cm,keepaspectratio]{monitoring-dashboard.pdf}

\vspace{0.1cm}
\footnotesize
\textbf{Not a pipeline --- a loop.} Data $\to$ Train $\to$ Deploy $\to$ Monitor $\to$ \textcolor{errorstroke}{\textbf{Retrain}}.\\
A/B testing in ML takes \textbf{4--6 weeks} (vs.\ 2--3 days for traditional software).

\end{frame}

% =============================================================================
\section{Lighthouses \& Pillars}
% =============================================================================

\begin{frame}{Lighthouse Models: Systems Detectives}
\note{
% -- LINK: The Iron Law showed that different workloads have different dominant
% terms. Lighthouse models are the concrete examples that make this tangible.
%
% -- NARRATE: ``Each lighthouse isolates a distinct bottleneck. ResNet-50 is
% compute-bound---dense matrix ops dominate. GPT-2/Llama is memory-bandwidth-
% bound---reads the full parameter set for each token. MobileNetV2 is
% latency-bound---must meet real-time deadlines on a phone. DLRM is data-I/O-
% bound---massive embedding tables with random access. KWS is power-bound---
% must run on a coin-cell battery for years. These five models recur in every
% chapter, testing concepts under different constraint regimes.''
%
% -- ENGAGE: ``Why do we need five lighthouse models, not one?'' Give 15
% seconds. Expected: because different models stress different parts of the
% system. One model cannot expose all bottleneck types.
%
% -- WARN: Students think ``one model to rule them all.'' Correct framing: the
% deployment spectrum spans $10^6\times$ in compute. No single model or hardware
% configuration works everywhere. That is why we need five detectives.
%
% -- FLEX: [CORE] Cast of characters for the entire course.
% IF SHORT: Show the table, name each lighthouse and its bottleneck, move on.
}

\footnotesize
\renewcommand{\arraystretch}{1.15}
\begin{tabular}{@{}lp{2cm}lp{3.5cm}@{}}
  \toprule
  \textbf{Lighthouse} & \textbf{Bottleneck} & \textbf{Iron Law} & \textbf{Question} \\
  \midrule
  \textbf{ResNet-50}    & \textcolor{computestroke}{Compute}  & $O/R_p$ & Fast arithmetic? \\
  \textbf{GPT-2/Llama}  & \textcolor{datastroke}{Memory-BW}  & $D/BW$  & Feed the accelerator? \\
  \textbf{MobileNetV2}  & \textcolor{routingstroke}{Latency}  & $L_{\text{lat}}$ & Meet deadlines? \\
  \textbf{DLRM}         & \textcolor{datastroke}{Data I/O}   & $D/BW$  & Move embeddings? \\
  \textbf{KWS}          & \textcolor{errorstroke}{Power}      & $P \times T$ & Coin-cell battery? \\
  \bottomrule
\end{tabular}

\vspace{0.1cm}
{\footnotesize Each lighthouse recurs in every chapter, testing concepts under different constraint regimes.}

\end{frame}

\begin{frame}{The Deployment Spectrum}
\note{
% -- LINK: Lighthouse models showed five different bottleneck types. The
% deployment spectrum shows the physical hardware that creates those bottlenecks.
%
% -- NARRATE: ``Nine orders of magnitude separate Cloud and TinyML. The H100
% has 80 GB of HBM3 and 989 TFLOPS. The Arduino Nicla has 384 KB of SRAM and
% 0.5 GFLOPS. That is a million-fold gap in compute. You cannot `shrink' a
% cloud model to TinyML---each tier requires fundamentally different engineering.
% Different algorithms, different precision, different memory management.''
%
% -- ENGAGE: ``Can you run GPT-2 (1.5B parameters, ~6 GB) on a mobile device
% with 8 GB RAM?'' Give 15 seconds.
% Expected: barely---model fills most of RAM, leaving nothing for the OS or
% app. Deepen: ``What about KV cache for generation?''
%
% -- WARN: Students think ``just quantize it and it runs everywhere.'' Correct
% framing: quantization helps (4--8$\times$) but the $10^6\times$ gap between
% Cloud and TinyML requires entirely different model architectures, not just
% lower precision on the same model.
%
% -- FLEX: [CORE] Physical reality that governs deployment decisions.
% IF SHORT: Show the table, state the $10^6\times$ gap, move on (60 seconds).
}

% --- Layout: IMAGE TOP + compact table ---
\centering
\safeimg[width=0.88\textwidth,height=3.5cm,keepaspectratio]{deployment-spectrum.pdf}

\vspace{0.05cm}
\scriptsize
\renewcommand{\arraystretch}{0.95}
\begin{tabular}{@{}llll@{}}
  \toprule
  \textbf{Tier} & \textbf{RAM} & \textbf{Compute} & \textbf{Power} \\
  \midrule
  Cloud (H100)       & 80 GB  & 989 TF   & 700 W \\
  Edge (Orin NX)     & 32 GB  & 100 TOP  & 15--40 W \\
  Mobile             & 8 GB   & 17 TOP   & 3--5 W \\
  TinyML (Nicla)     & 384 KB & 0.5 GF   & 100 mW \\
  \bottomrule
\end{tabular}
\hfill
{\scriptsize\textcolor{errorstroke}{\textbf{Gap:}} $10^{5}\!\times$ memory, $10^{6}\!\times$ compute}

\end{frame}

\begin{frame}{Amdahl's Law: A Systems Pitfall}
\note{
% -- LINK: The deployment spectrum showed different hardware tiers. Amdahl's
% Law warns that optimizing one component of a pipeline has bounded returns.
%
% -- NARRATE: ``Even a 3$\times$ speedup on inference yields only 23\%
% end-to-end improvement when preprocessing dominates. Look at the table:
% inference drops from 45 ms to 15 ms---a dramatic improvement. But
% preprocessing is still 60 ms and postprocessing is 25 ms. Total goes from
% 130 ms to 100 ms. That is only 23\% faster, not 3$\times$. This is why
% you must profile the whole pipeline before optimizing any component.''
%
% -- ENGAGE: ``If you could magically make inference take 0 ms, what would
% the total pipeline time be?'' Give 10 seconds.
% Expected: 60 + 0 + 25 = 85 ms. Maximum speedup = 130/85 = 1.53$\times$.
% Amdahl's Law: the non-accelerated fraction (65\%) sets the floor.
%
% -- WARN: Students focus on model inference speed and ignore the rest of the
% pipeline. Correct framing: a 3$\times$ speedup on 35\% of the pipeline
% yields only 23\% end-to-end. Profile the whole pipeline first.
%
% -- FLEX: [OPTIONAL] Important concept but can be abbreviated.
% IF SHORT: State the punchline: ``3$\times$ on one stage = 23\% end-to-end.
% Profile first.'' (30 seconds)
}

% --- Layout: IMAGE TOP + compact table below ---
\centering
\safeimg[width=0.85\textwidth,height=2.8cm,keepaspectratio]{amdahls-pipeline.pdf}

\vspace{0.05cm}
\scriptsize
\renewcommand{\arraystretch}{0.95}
\begin{tabular}{@{}lrrr@{}}
  \toprule
  \textbf{Stage} & \textbf{Before} & \textbf{After} & \textbf{Change} \\
  \midrule
  Preprocessing  & 60 ms & 60 ms & --- \\
  Inference      & 45 ms & 15 ms & \textcolor{datastroke}{3$\times$ faster} \\
  Postprocessing & 25 ms & 25 ms & --- \\
  \midrule
  \textbf{Total} & \textbf{130 ms} & \textbf{100 ms} & \textcolor{errorstroke}{\textbf{only 23\%}} \\
  \bottomrule
\end{tabular}
\par\vspace{0.05cm}
{\footnotesize\textcolor{errorstroke}{\textbf{Amdahl's Law:}} 3$\times$ speedup on 35\% of the pipeline $\to$ only 23\% end-to-end.}

\end{frame}

% --- ACTIVE LEARNING 3: Discussion ---
\begin{frame}{Discussion: Which Pillar Fails First?}
\note{
% -- LINK: Students have learned D-A-M, the Iron Law, lighthouse models, and
% Amdahl's Law. This discussion asks them to synthesize: which engineering
% discipline fails when a system breaks in production?
%
% -- NARRATE: ``A startup deploys a highly accurate model. Within 3 months,
% customer complaints spike. Which engineering discipline failed?'' Point to
% the five options. ``There is no single right answer---the point is that all
% pillars are interdependent. A monitoring gap lets drift go undetected. A
% data engineering gap means the training data was not representative. A
% deployment gap means the serving infrastructure could not handle scale.''
%
% -- ENGAGE: Turn-and-talk for 90 seconds. Cold-call 2--3 pairs.
% Accept any well-reasoned answer. Push students to name a specific mechanism:
% not ``monitoring failed'' but ``monitoring failed to detect distribution
% drift in the input features.''
%
% -- WARN: Students will pick one answer and defend it. Correct framing:
% multiple pillars can fail simultaneously, and the root cause is often
% upstream of where the symptom appears.
%
% -- FLEX: [CORE] Synthesis exercise for the chapter.
% IF SHORT: Do a show-of-hands poll instead of turn-and-talk.
}

\centering
\vspace{0.8cm}
{\large\bfseries Turn and Talk \textcolor{midgray}{(90 seconds)}}

\vspace{0.6cm}
{\large A startup deploys a highly accurate model.\\
Within 3 months, customer complaints spike.\\[0.3cm]
\alert{Which engineering discipline failed?}}

\vspace{0.6cm}
\begin{columns}[c]
  \begin{column}{0.18\textwidth}\centering\small Data\\ Eng.\end{column}
  \begin{column}{0.18\textwidth}\centering\small Training\\ Systems\end{column}
  \begin{column}{0.18\textwidth}\centering\small Deploy\\ Infra\end{column}
  \begin{column}{0.18\textwidth}\centering\small Ops \&\\ Monitor\end{column}
  \begin{column}{0.18\textwidth}\centering\small Ethics \&\\ Gov.\end{column}
\end{columns}

\end{frame}

\begin{frame}{Case Study: Voice Assistant End-to-End}
\note{
% -- LINK: Students have seen the deployment spectrum and lighthouse models
% in isolation. This slide shows a real system that spans multiple tiers.
%
% -- NARRATE: ``A voice assistant is three ML systems stitched together.
% Wake-word detection runs on TinyML at under 1 mW---always on, listening.
% Speech-to-text runs on the mobile SoC at 3--5 W---on-device for privacy.
% Language understanding runs in the cloud at 700 W---because the model is
% too large for the device. Total latency budget: 200 ms. KWS takes 10 ms.
% STT takes 80 ms. Cloud round-trip plus LLM: 110 ms. Every millisecond is
% accounted for. This is AI Engineering: three D-A-M analyses, three Iron
% Law calculations, one user experience.''
%
% -- ENGAGE: ``Which stage is most likely to miss the 200 ms budget?'' Give
% 15 seconds. Expected: cloud round-trip (network latency is variable).
%
% -- WARN: Students think of this as ``one system.'' Correct framing: it is
% three systems with three different constraint regimes, stitched together.
%
% -- FLEX: [CORE] Concrete deployment case study.
% IF SHORT: Show the pipeline, state the 200 ms budget, move on.
}

\scriptsize
\renewcommand{\arraystretch}{1.1}
\begin{tabular}{@{}llllll@{}}
  \toprule
  \textbf{Stage} & \textbf{Model} & \textbf{Tier} & \textbf{Power} & \textbf{Latency} & \textbf{Bottleneck} \\
  \midrule
  Wake-word   & KWS DS-CNN    & TinyML & $<$1 mW  & 10 ms  & \textcolor{errorstroke}{Power} \\
  Speech-to-text & Whisper-tiny & Mobile & 3--5 W  & 80 ms  & \textcolor{computestroke}{Compute} \\
  Language    & LLM (cloud)    & Cloud  & 700 W   & 110 ms & \textcolor{datastroke}{Memory BW} \\
  \midrule
  \textbf{Total} & & & & \textbf{200 ms} & \\
  \bottomrule
\end{tabular}

\vspace{0.15cm}
\small
\begin{columns}[T]
  \begin{column}{0.48\textwidth}
    \begin{mlsyscard}{crimson}
    \textbf{Three D-A-M analyses}\\
    {\footnotesize Three Iron Law calculations\\Three constraint regimes\\One user experience: 200 ms}
    \end{mlsyscard}
  \end{column}
  \begin{column}{0.48\textwidth}
    \begin{mlsyscard}{routingstroke}
    \textbf{Network is the wildcard}\\
    {\footnotesize Cloud round-trip: 20--100 ms variable. If network spikes, the entire 200 ms budget breaks.}
    \end{mlsyscard}
  \end{column}
\end{columns}

\end{frame}

\begin{frame}{The Five Pillars of ML Systems}
\note{
% -- LINK: The discussion surfaced multiple engineering disciplines that can
% fail. This slide names the five pillars and maps them to the course.
%
% -- NARRATE: ``Five pillars, each essential. Data Engineering is the fuel---
% garbage in, garbage out. Training Systems is the engine---compute, memory,
% optimization. Deployment Infrastructure is the factory---serving, scaling,
% latency. Ops and Monitoring is the cockpit---drift detection, alerting,
% retraining triggers. Ethics and Governance is the guardrails---fairness,
% privacy, compliance. Teams lacking expertise in any single pillar face the
% 60--85\% failure rate we discussed in the Course Overview.''
%
% -- ENGAGE: ``Which pillar do you think gets the least attention in
% industry?'' Give 10 seconds. Show of hands.
% Expected: Ethics/Governance or Ops/Monitoring. Both are chronically
% under-invested relative to Training Systems.
%
% -- WARN: Students think strong ML expertise compensates for weak systems.
% Correct framing: teams with strong ML but weak systems achieve only 10--20\%
% of throughput targets. All five pillars must be staffed.
%
% -- FLEX: [CORE] Organizing framework for the rest of the course.
% IF SHORT: Name the five pillars, state the 60--85\% failure rate, move on.
}

\footnotesize
\begin{columns}[T]
  \begin{column}{0.42\textwidth}
    \safeimg{book_pillars.png}
  \end{column}
  \begin{column}{0.55\textwidth}
    \begin{enumerate}\setlength\itemsep{0pt}
      \item \textcolor{datastroke}{\textbf{Data Engineering}} --- fuel
      \item \textcolor{computestroke}{\textbf{Training Systems}} --- engine
      \item \textcolor{routingstroke}{\textbf{Deployment Infra}} --- factory
      \item \textcolor{crimson}{\textbf{Ops \& Monitoring}} --- cockpit
      \item \textcolor{errorstroke}{\textbf{Ethics \& Governance}} --- guardrails
    \end{enumerate}

    \vspace{0.1cm}
    \begin{mlsyscard}{errorstroke}
    \textbf{60--85\%} of ML projects fail to reach production --- systems gaps, not algorithms.
    \end{mlsyscard}
  \end{column}
\end{columns}
\mlsyscite{Paleyes et al., ACM Computing Surveys 2022}

\end{frame}

\begin{frame}{Worked Example: GPT-2 Through Five Pillars}
\note{
% -- LINK: The Five Pillars named abstract disciplines. This slide maps GPT-2
% through all five to show they are concrete, not theoretical.
%
% -- NARRATE: ``GPT-2 is 1.5B parameters, about 6 GB at FP16. Walk through
% each pillar. Data Engineering: 40 GB WebText, 8M documents, cleaned and
% deduplicated. Training Systems: 256 TPU-v3 cores, 1 week, compute-bound.
% Deployment: memory-bandwidth-bound at inference---reads full 6 GB per token.
% Ops: monitor perplexity drift as language evolves. Ethics: training data
% contains biases from the internet. Every pillar is load-bearing.''
%
% -- ENGAGE: ``Which pillar is most likely to cause a production failure for
% GPT-2?'' Give 15 seconds. Expected: Ops (drift) or Ethics (bias).
%
% -- WARN: Students think training is the hard part. Correct framing: training
% is a one-time cost; deployment, ops, and ethics are ongoing.
%
% -- FLEX: [OPTIONAL] Can compress to 60 seconds if running short.
% IF SHORT: Walk through the table without discussion.
}

\scriptsize
\renewcommand{\arraystretch}{1.1}
\textbf{GPT-2 (1.5B params, 6 GB FP16) through all five pillars:}

\vspace{0.1cm}
\begin{tabular}{@{}lll@{}}
  \toprule
  \textbf{Pillar} & \textbf{GPT-2 Specifics} & \textbf{Key Metric} \\
  \midrule
  \textcolor{datastroke}{\textbf{Data Eng.}}      & 40 GB WebText, 8M docs, dedup'd & Tokens/doc quality \\
  \textcolor{computestroke}{\textbf{Training}}     & 256 TPU-v3, 1 week & 6$\times$N$\times$D FLOPs \\
  \textcolor{routingstroke}{\textbf{Deployment}}   & Memory-BW-bound: reads 6 GB/token & $D_{\text{vol}}/BW$ dominates \\
  \textcolor{crimson}{\textbf{Ops}}                & Monitor perplexity as language evolves & Degradation Eq. \\
  \textcolor{errorstroke}{\textbf{Ethics}}         & Internet biases in training data & Toxicity rate \\
  \bottomrule
\end{tabular}

\vspace{0.15cm}
\small
\begin{mlsyscard}{crimson}
{\footnotesize Training is a one-time cost. Deployment, ops, and ethics are \textbf{ongoing}. The five pillars are not sequential stages---they are concurrent disciplines.}
\end{mlsyscard}

\end{frame}

% =============================================================================
\section{Wrap-Up}
% =============================================================================

\begin{frame}{Fallacies}
\note{
% -- LINK: The chapter covered D-A-M, Iron Law, Degradation Equation,
% lighthouses, and five pillars. These fallacies test whether students
% internalized the key distinctions.
%
% -- NARRATE: ``Four misconceptions, each with quantitative evidence. First:
% better algorithms do not automatically produce better systems---ViT is 1--2\%
% more accurate but needs 4$\times$ the FLOPs, violating latency SLAs. Second:
% high benchmark accuracy does not mean production readiness---94\% on
% benchmark drops to 78--82\% due to distribution shift. Third: deploy-once
% does not work---rec systems drop 85\% to 79.2\% in 6 months. Fourth:
% testing cannot verify ML systems---the input space has $10^{362,000}$
% configurations.''
%
% -- FLEX: [CORE] Fallacies are essential for correcting misconceptions.
% IF SHORT: Cover just the first two fallacies.
}

\small
\textbf{Fallacy:} \textit{Better algorithms automatically produce better systems.}\\
{\footnotesize ViT: 1--2\% higher accuracy than ResNet but 4$\times$ FLOPs. A model violating latency requirements has zero utility.}

\vspace{0.1cm}
\textbf{Fallacy:} \textit{High benchmark accuracy means production readiness.}\\
{\footnotesize Sentiment model: 94\% on benchmark $\to$ 78--82\% in production due to distribution shift.}

\vspace{0.1cm}
\textbf{Fallacy:} \textit{Deploy once and let it run.}\\
{\footnotesize Rec.\ system: 85\% $\to$ 79.2\% in 6 months. Fraud detection: 5--10\% decay per quarter.}

\vspace{0.1cm}
\textbf{Fallacy:} \textit{ML systems can be verified by testing.}\\
{\footnotesize Input space: $256^{150{,}528}$ configs. Test set: 50{,}000. The gap is effectively infinite.}

\end{frame}

\begin{frame}{Pitfalls}
\note{
% -- LINK: Fallacies named misconceptions. Pitfalls name operational mistakes
% students will make in practice.
%
% -- NARRATE: ``Three pitfalls from real systems. First: optimizing components
% without system context---Amdahl's Law means a 3$\times$ inference speedup
% yields only 23\% end-to-end when preprocessing dominates. Second: treating ML
% as traditional software---ML degrades 10--15\% over 3--6 months silently;
% unit tests are insufficient. Third: assuming ML expertise alone is
% sufficient---teams with strong ML but weak systems achieve only 10--20\% of
% throughput targets.''
%
% -- FLEX: [CORE] Operational lessons.
% IF SHORT: Cover just the first two pitfalls.
}

\footnotesize
\textbf{Pitfall:} \textit{Optimizing components without considering system interactions.}\\
A 3$\times$ inference speedup (45$\to$15 ms) yields only \textbf{23\%} end-to-end improvement when preprocessing dominates. \textbf{Amdahl's Law governs the whole pipeline.}

\vspace{0.15cm}
\textbf{Pitfall:} \textit{Treating ML systems as traditional software.}\\
ML systems degrade 10--15\% over 3--6 months silently. Unit tests are insufficient; \textbf{monitoring infrastructure} is required from day one.

\vspace{0.15cm}
\textbf{Pitfall:} \textit{Assuming ML expertise alone is sufficient.}\\
Teams with strong ML but weak systems achieve only \textbf{10--20\%} of throughput targets. The five disciplines require integrated expertise.

\end{frame}


\begin{frame}{Muddiest Point}
\note{
% -- NARRATE: ``One-minute paper. Write down the one concept from today that
% was most confusing or unclear. Hand it in on your way out or submit
% digitally. I will open the next lecture by addressing the top 2--3 muddiest
% points.''
%
% -- FLEX: [CORE] Diagnostic tool (Angelo \& Cross, 1993). Never skip.
% Responses shape the opening of the next lecture.
}

\centering
\Large\bfseries One-minute paper\\[0.5cm]
\normalsize Write down the \textbf{muddiest point} from today's lecture.\\[0.2cm]
{\small What concept was most confusing or unclear?}\\[0.5cm]
{\footnotesize\textcolor{midgray}{(Hand in on your way out --- or submit digitally)}}
\end{frame}

% --- RETRIEVAL PRACTICE ---
\begin{frame}{What Were the Key Ideas?}
\note{
% -- NARRATE: ``Close your notes. Write down the 4 most important concepts
% from today. 90 seconds, no peeking. This is retrieval practice---recalling
% information strengthens memory more than re-reading it.''
% Walk around the room. Do NOT show the next slide yet.
%
% -- FLEX: [CORE] Retrieval practice (Roediger \& Karpicke, 2006).
% Never skip. Even 60 seconds of recall is more valuable than 5 minutes
% of re-reading.
}

\centering
\vspace{1.5cm}
{\Large\bfseries Close your notes.}

\vspace{0.8cm}
{\large Write down the \textbf{4 most important concepts} from today.}

\vspace{0.8cm}
{\normalsize\textcolor{midgray}{90 seconds --- no peeking.}}

\end{frame}

\begin{frame}{Key Takeaways}
\note{
% -- NARRATE: ``Reveal time. Walk through each bullet. D-A-M: three
% interdependent axes---optimizing one shifts the bottleneck. Silent failure:
% the Degradation Equation quantifies drift. Iron Law: three terms, slowest
% dominates. Bitter Lesson: scale wins, but 44$\times$ algorithmic efficiency
% compounds too. AI Engineering: jointly optimizing all axes under shifting
% distributions. Five pillars: missing any one leads to the 60--85\% failure
% rate. Deployment spectrum: $10^6\times$ gap means different engineering per
% tier. Emphasize the quantitative anchors: 23\% Amdahl's, 362,000-digit
% verification gap, 44$\times$ efficiency, 60--85\% failure rate.''
%
% -- FLEX: [CORE] Summary slide. Walk through all bullets.
% IF SHORT: Hit just the three frameworks (D-A-M, Iron Law, Degradation).
}

\footnotesize
\begin{itemize}\setlength\itemsep{1pt}
  \item \textbf{\DAM{} taxonomy}: Data, Algorithm, Machine are interdependent. Optimizing one shifts the bottleneck.
  \item \textbf{Silent failure}: Degradation Equation quantifies accuracy decay. Set retraining triggers from drift.
  \item \textbf{Iron Law}: $T = D_{\text{vol}}/BW + O/(R_{\text{peak}} \cdot \eta) + L_{\text{lat}}$. Slowest term dominates.
  \item \textbf{Bitter Lesson}: Scale wins --- but algorithmic efficiency (44$\times$) compounds alongside it.
  \item \textbf{AI Engineering}: Jointly optimizing all \DAM{} axes under continuously shifting distributions.
  \item \textbf{Five pillars}: Data, Training, Deployment, Ops, Ethics. Missing any $\to$ 60--85\% failure.
  \item \textbf{Deployment spectrum}: Cloud to TinyML spans $10^6\!\times$ compute. Each tier needs different \DAM{} balance.
\end{itemize}

\end{frame}

\begin{frame}{References}
\note{
% -- NARRATE: ``Six canonical papers. Sculley for hidden technical debt.
% Sutton for the Bitter Lesson. Karpathy for Software 2.0. Hernandez and
% Brown for algorithmic efficiency. Patterson and Hennessy for the quantitative
% approach. Paleyes for deployment challenges.''
%
% -- FLEX: [OPTIONAL] Show briefly. Students can look up papers on their own.
}

\small
\mlsysref{Sculley+15}{Sculley et al. ``Hidden Technical Debt in ML Systems.'' NeurIPS 2015.}
\mlsysref{Sutton19}{R. Sutton. ``The Bitter Lesson.'' 2019.}
\mlsysref{Karpathy17}{A. Karpathy. ``Software 2.0.'' 2017.}
\mlsysref{Hernandez+20}{Hernandez \& Brown. ``Measuring Algorithmic Efficiency.'' 2020.}
\mlsysref{P\&H21}{Patterson \& Hennessy. \textit{Comp.\ Arch.: A Quantitative Approach}. 6th ed.}
\mlsysref{Paleyes+22}{Paleyes et al. ``Challenges in Deploying ML.'' ACM Comp.\ Surveys, 2022.}

\end{frame}

\begin{frame}{Next Lecture: ML Systems in the Wild}
\note{
% -- LINK: This chapter established three frameworks. The next chapter asks:
% where should a model actually run?
%
% -- NARRATE: ``Physical laws dictate where a model can run. Speed of light
% makes distant cloud servers useless for emergency braking---36 ms minimum
% round trip from CA to VA, before any computation. The Power Wall limits
% mobile to 3--5 W. The Memory Wall limits TinyML to 384 KB. The question
% `where should a model run?' is the entire next chapter.''
%
% -- FLEX: [CORE] Forward hook. Assign pre-reading if applicable.
% IF SHORT: ``Next class: where should a model run? Read Chapter 2.''
}

\footnotesize
\begin{columns}[c]
  \begin{column}{0.30\textwidth}
    \centering
    \safeimg[height=2.5cm]{nvidia-h100.png}\\
    \textcolor{computestroke}{\textbf{Cloud}} --- 700 W
  \end{column}
  \begin{column}{0.30\textwidth}
    \centering
    \safeimg[height=2.5cm]{jetson-orin-nx.jpg}\\
    \textcolor{routingstroke}{\textbf{Edge}} --- 15--40 W
  \end{column}
  \begin{column}{0.30\textwidth}
    \centering
    \safeimg[height=2.5cm]{arduino-nicla-vision.jpg}\\
    \textcolor{datastroke}{\textbf{TinyML}} --- 100 mW
  \end{column}
\end{columns}

\vspace{0.15cm}
\centering
{\normalsize Where should an ML model run?}\\[0.1cm]
{\small\textbf{The speed of light, thermodynamics, and battery physics decide.}}

\end{frame}


% =============================================================================
% BACKUP SLIDES
% =============================================================================
\appendix

\begin{frame}{Backup: ResNet-50 on H100 --- Compute or Memory?}
\note{
% -- NARRATE: Use as an additional exercise if students finish the main
% exercise early. ResNet-50: 8 GFLOPs per image. H100: 989 TFLOPS peak,
% 40\% utilization. Inference time = 8e9 / (989e12 * 0.4) = 20 microseconds.
% But actual time is ~1 ms. Why? Memory bandwidth: loading 97.5 MB of
% parameters at 3,350 GB/s = 29 microseconds. Plus framework overhead.
% The model is memory-bound at batch=1, not compute-bound.
%
% -- FLEX: [OPTIONAL] Additional quantitative exercise.
}
\small
\textbf{ResNet-50 inference on H100 (batch=1):}

\vspace{0.15cm}
{\footnotesize
\renewcommand{\arraystretch}{1.15}
\begin{tabular}{@{}llll@{}}
  \toprule
  \textbf{Term} & \textbf{Variables} & \textbf{Calculation} & \textbf{Time} \\
  \midrule
  \textcolor{computestroke}{Compute} & 8 GFLOPs, 989 TF, $\eta$=0.4 & $8\text{G}/(396\text{T})$ & \textbf{20 $\mu$s} \\
  \textcolor{datastroke}{Data}    & 97.5 MB params, 3,350 GB/s & $97.5\text{M}/3{,}350\text{G}$ & \textbf{29 $\mu$s} \\
  \textcolor{errorstroke}{Overhead} & Framework, kernel launch & --- & \textbf{$\sim$950 $\mu$s} \\
  \midrule
  & \textbf{Actual measured} & & \textbf{$\sim$1 ms} \\
  \bottomrule
\end{tabular}
}

\vspace{0.15cm}
\textbf{Surprise:} Neither compute nor data dominates---\textcolor{errorstroke}{overhead} does at batch=1. At batch=256, compute catches up and becomes the bottleneck.
\end{frame}

\begin{frame}{Backup: Degradation Equation Worked Example}
\note{
% -- NARRATE: Use if students struggle with the Degradation Equation.
% Walk through month by month: $A(t) = 0.85 - 0.012 \times \Delta(t)$.
% At month 6, $\Delta = 0.48$, so $A = 0.85 - 0.012 \times 0.48 / 0.012
% = 79.2\%$. Set retraining trigger at $A(t) < 0.80$ (month $\sim$5).
%
% -- FLEX: [OPTIONAL] Backup slide. Use only if students need scaffolding.
}
\small
Use if students struggle with the Degradation Equation.

\small
\textbf{Scenario:} Recommendation system, $A_0 = 0.85$, $\alpha = 0.012$/month.

\vspace{0.2cm}
{\footnotesize
\begin{tabular}{@{}lrrr@{}}
  \toprule
  \textbf{Month} & $\Delta(P)$ & $A(t)$ & \textbf{Status} \\
  \midrule
  0 & 0.00 & 85.0\% & Deployed \\
  3 & 0.15 & 83.2\% & OK \\
  6 & 0.48 & 79.2\% & Alert threshold \\
  12 & 1.10 & 71.8\% & SLA violation \\
  \bottomrule
\end{tabular}
}

\vspace{0.15cm}
\textbf{Key:} Set retraining trigger at $A(t) < 0.80$ (month $\sim$5).
\end{frame}

\begin{frame}{Backup: Amdahl's Law --- Why System Thinking Matters}
\note{
% -- NARRATE: Use if students question the 23\% end-to-end result.
% Walk through the formula: $\text{Speedup} = 1 / ((1-p) + p/S)$.
% With $p = 0.35$ and $S = 3$: $1 / (0.65 + 0.117) = 1.30$.
% Even with infinite speedup on inference: $1 / 0.65 = 1.54$.
%
% -- FLEX: [OPTIONAL] Backup slide. Use only if students need the math.
}
\small
Use if students question the 23\% end-to-end result.

\small
$$\text{Speedup} = \frac{1}{(1-p) + p/S}$$

\vspace{0.1cm}
If inference is 35\% of the pipeline and you speed it up 3$\times$:
$$\text{Speedup} = \frac{1}{0.65 + 0.35/3} = \frac{1}{0.767} = 1.30$$

\vspace{0.1cm}
Only \textbf{30\% improvement} despite 3$\times$ faster inference. The other 65\% (preprocessing + postprocessing) is untouched. \textbf{Profile the whole pipeline first.}
\end{frame}

\end{document}
